\documentclass[11pt]{article}

% -----------------------------------------------------------------------
% 中文版：本地阅读用 [final]，投稿时改为 [review]
% -----------------------------------------------------------------------
\usepackage[final]{acl}

\usepackage{xeCJK}
\setCJKmainfont{Songti SC}             % macOS 宋体

\usepackage{latexsym}
\usepackage{microtype}
\usepackage{graphicx}
\usepackage{amsmath}
\usepackage{amssymb}
\usepackage{booktabs}
\usepackage{multirow}
\usepackage{xcolor}
\usepackage{url}

\renewcommand\outauthor{%
    \begin{tabular}[t]{c}
    \bfseries 匿名 EMNLP 2026 投稿
    \end{tabular}}

\newcommand{\method}{RDT}
\newcommand{\methodfull}{拒绝方向移植}

\title{拒绝方向移植（RDT/RDT+）：\\
  将文本主干安全对齐迁移至\\
  视觉-语言-动作模型的动作 Token 通路}

\author{匿名 EMNLP 2026 投稿}

\begin{document}
\maketitle

\begin{abstract}
视觉-语言-动作（VLA）模型（如 OpenVLA）通过覆写 LLM 词表的末尾来输出离散动作 token，从而复用预训练语言模型。我们认为，这种 tokenization 方案制造了一个\textit{结构性安全缺口}：从 chat 对齐 LLM 继承的拒绝回路在几何上与 base 主干 VLA 的动作 token 通路解耦，导致有害指令直接转化为不安全的物理动作而不触发任何拒绝。我们提出\methodfull{}（RDT/RDT+），一组无需训练的推理时干预方法，从与目标 VLA 共享预训练初始化的外部对齐兄弟 LLM（Llama-2-chat）中提取拒绝方向，并在动作生成位置重新注入。基础 RDT 针对 KV 缓存中的动作 token 位置；RDT+ 扩展为生成阶段感知的双阶段注入，同时引导解码第一个动作 token 的指令表示，实现对完整 7 token 动作序列的全覆盖。我们通过方向来源对照实验——将拒绝方向与等范数随机单位向量对比——来验证移植方向的特定几何性质（而非仅仅是幅度）是有效成分。在 \textsc{LIBERO-Long}、\textsc{SafeAgentBench} 和指令条件化 \textsc{AdvBench} 分集上，我们在文本越狱和通用对抗补丁攻击下评测 RDT+，与提示层、隐层投影和 LoRA 微调防御对比，并在白盒自适应攻击者下进行压力测试。据我们所知，本研究是首次在不重训练底层模型的情况下，将文本对齐视为可移植资产注入非文本输出通路的工作。
\end{abstract}

\section{引言}
\label{sec:intro}

大型多模态模型从语言模型主干继承其安全对齐能力，而该主干的 RLHF 训练完全在文本上进行。每一种额外接入该主干的输入或输出模态——图像、音频、视频，或如本文所研究的 7-DoF 机器人动作 token——都会开辟一条拒绝回路从未见过的通路。当对齐能力迁移成功时，我们获得更安全的多模态助手；当迁移失败时，轻则出现多模态幻觉，重则导致物理危害。

视觉-语言-动作（VLA）模型处于这一谱系中危害最大的一端。一个不安全的动作并非一个不安全的\emph{句子}：它是机械臂夹住孩子的手、将液体倒入明火、将武器对准一个人。近期对抗性研究表明，目前所有公开研究的 VLA 模型均可被通过简单的语言指令或微小的图像扰动诱导生成此类动作 \cite{wu2024exploringadversarial,badrobot2025,shawshank2025}，而现有 VLA 防御要么在安全强化学习框架下对整个策略进行端到端重训练 \cite{safevla2025}，要么依赖于不包含语义危害概念的几何碰撞规避方法 \cite{vlsa2025}。

\paragraph{动作 token VLA 中的结构性安全缺口。}
我们观察到，事实上的开源 VLA——\textsc{OpenVLA} \cite{openvla2024}——建立在 Llama-2-7b-\emph{base} 之上，这是一个\emph{从未经过 RLHF 对齐}的模型，并通过覆写 Llama 词表的最后 256 个 token 来编码动作词汇。这两个设计选择是当前社区的惯例；我们指出，二者叠加会产生一个现有多模态安全防御方法无法跨越的几何缺口。首先，主干本身没有拒绝方向：它从未被训练过拒绝。其次，动作 token 在经历大量行为克隆微调后，驻留在残差流的一个子空间中，该子空间与若主干经过 chat 对齐后本应存在的"有害/无害"自然语言轴近乎正交。我们在 \S\ref{sec:method:diagnostic} 中对该缺口进行定量分析：在已对齐 LVLM 上能够清晰区分有害与无害样本的安全层探针 \cite{securitytensors2025}，在 \textsc{OpenVLA} 的动作解码通路上信号显著减弱。

\paragraph{跨模型拒绝方向移植。}
我们提出一种无需重训练 VLA 的修复方案。\methodfull{}（\method{}）从 Llama-2-7b-\emph{chat} 中提取拒绝方向——该模型与 \textsc{OpenVLA} 主干共享预训练初始化但已获得 RLHF 对齐——并在推理时仅在 \textsc{OpenVLA} 的动作 token 位置将其重新注入。该方法以单个前向钩子实现，可与现有提示层防御组合使用。

两项先前工作在方法上与我们相邻，但均不足以处理本文场景。\textsc{VLM-Guard} \cite{vlmguard2025} 提取了类似的引导方向，但将其施加于所有 token，过度扰动了良性的视觉-语言处理。\textsc{SARSteer} \cite{sarsteer2025} 在音频语言模型中引导文本衍生的拒绝方向，但关键区别在于其从目标模型自身提取方向，而当目标主干（Llama-2-base）完全没有拒绝响应时，该方法零样本不适用。

\paragraph{贡献。}
\begin{itemize}
    \item 我们识别出动作 token VLA 特有的结构性安全缺口，该缺口源于\emph{未对齐}主干与学习到的离散动作词汇的组合（\S\ref{sec:method:diagnostic}）。
    \item 我们提出 \method{} 与 \method{}+，两种无训练的推理时方法，能将拒绝方向从外部对齐 LLM 移植到 VLA 中。基础 \method{} 以动作 token 位置为目标；\method{}+ 增加了覆盖全部 7 个动作 token（含 $a_1$）的生成阶段感知双阶段注入（\S\ref{sec:method:rdt}--\ref{sec:method:dualphase}），并提供 rank-$k$ 子空间变体（\S\ref{sec:method:rankk}）和内容自适应缩放头（\S\ref{sec:method:adaptive}）。
    \item 我们在三个评测集（\textsc{LIBERO-Long}、\textsc{SafeAgentBench}、指令条件化的 \textsc{AdvBench} 分集）和四类攻击家族（文本越狱、\textsc{UADA}、\textsc{UPA}、\textsc{TMA} 对抗补丁）上设计了系统评测，将 \method{}+ 与提示层、隐藏状态投影和 LoRA 微调基线进行对比（\S\ref{sec:exp}）。
    \item 我们引入方向来源对照实验——拒绝方向 vs.\ 等范数随机单位向量——以确认跨模型移植的拒绝方向几何特性（而非仅仅是扰动量级）是有效成分（\S\ref{sec:exp:ablation}）。
    \item 我们在白盒自适应攻击者下对 \method{}+ 进行压力测试，并分析动作空间中"拒绝"的含义，包括安全层探针恢复分析（\S\ref{sec:exp:adaptive}、\S\ref{sec:analysis:action-space}、\S\ref{sec:analysis:probe-recovery}）。
\end{itemize}

\section{相关工作}
\label{sec:related}

\paragraph{多模态安全对齐及其脆弱性。}
研究表明，视觉指令微调会系统性地侵蚀底层 LLM 的安全对齐 \cite{vlguard2024}，后续分析 \cite{safemirage2025} 指出，SFT 获得的拒绝能力在很大程度上是表浅的文本相关性。多模态偏好对微调 \cite{spavl2024} 和跨模态标题引导对齐 \cite{cmrt2025} 从训练阶段入手解决这一问题。推理时提示防御方法包括 \textsc{AdaShield} \cite{adashield2024} 和 \textsc{ECSO} \cite{ecso2024}，而 \textsc{MMJ-Bench} \cite{mmjbench2024} 表明单一防御无法覆盖所有攻击类型。

\paragraph{表示层安全干预。}
\citet{arditi2024refusal} 表明，在 chat 对齐的 LLM 中，拒绝行为由残差流中的单一方向所介导（通过有害/无害提示的均值差分提取）。关键地，后续工作证实，将 chat 模型的拒绝方向移植到\emph{相同预训练}的 base LLM 中，可将拒绝率从接近零提升至 88\% 以上 \cite{baselmsrefuse2024}，直接为我们所依赖的跨模型几何假设提供了实证基础。\textsc{Circuit Breakers} \cite{circuitbreakers2024} 通过 LoRA 适配器重定向有害表示。在 VLM 中，\textsc{Security Tensors} \cite{securitytensors2025} 识别出"安全层"的连续范围（LLaMA-3.2-11B-Vision 中为第 9--20 层），这些层的激活对有害文本输入响应，而对有害图像输入沉默；我们在 \textsc{OpenVLA} 的动作解码通路上复现了该探针。\textsc{CMRM} \cite{cmrm2025} 提供测试时位移向量；\textsc{VLM-Guard} \cite{vlmguard2025} 通过 SVD 从对齐 LLM 中提取 rank-$k$ 引导方向，并在推理时将其施加于 VLM 的最终输入 token 位置。\textsc{ASTRA} \cite{astra2025} 从视觉对抗样本中推导引导方向。\textsc{SafeSteer} \cite{safesteer2025} 将 SVD 子空间投影应用于 VLM 的通用越狱防御。

与我们工作最相近的是 \textsc{SARSteer} \cite{sarsteer2025}，该同期工作在音频语言模型中引导文本衍生的拒绝方向。关键区别在于：\textsc{SARSteer} 从目标模型自身提取方向，其主干已具备拒绝响应。而本文目标 VLA 建立在 \emph{base} LLM 之上，完全没有拒绝响应可供提取——在此场景下 \textsc{SARSteer} 的方法零样本不可用。我们从外部对齐的兄弟模型中提取方向并跨模型移植。

\paragraph{VLA 的激活引导用于行为控制。}
独立于安全问题，近期可解释性工作表明，激活引导对于 VLA 的\emph{行为}控制是有效的：\citet{mechinterpvla2025} 展示了在 \textsc{OpenVLA} 和 \textsc{$\pi_0$} 的特定 FFN 位置添加学习到的引导向量，可以可靠地修改运动速度和轨迹曲率。这为我们的安全干预建立了前提——\textsc{OpenVLA} 的残差流具有可被方向引导的几何结构——尽管其目标是行为控制而非危害拒绝，方向来源是分布内 VLA 数据而非跨模型拒绝方向，且未限制注入到动作 token 位置。

\paragraph{VLA 模型及其安全性。}
\textsc{OpenVLA} \cite{openvla2024} 和 \textsc{OpenVLA-OFT} \cite{openvlaoft2025} 代表了两种主流动作解码器：覆写 Llama 词表的离散动作 token，以及连续 L1 回归头。\textsc{$\pi_0$} \cite{pi02024} 在连续动作上使用流匹配。攻击性工作 \cite{wu2024exploringadversarial,badrobot2025,shawshank2025} 揭示了一个一致的结构性观察：VLA 没有拒绝训练的类比，且通用对抗补丁可跨任务迁移。防御方面，\textsc{SafeVLA} \cite{safevla2025} 将问题构建为约束马尔可夫决策过程，在安全强化学习下微调策略；\textsc{VLSA}/\textsc{AEGIS} \cite{vlsa2025} 在动作输出层施加控制论碰撞规避。这两种方法均未在表示层进行干预，即都没有询问 LLM 主干在生成动作之前是否本可感知该请求是有害的。

\paragraph{我们的定位。}
表~\ref{tab:positioning} 沿四个维度总结了设计空间：干预目标、阶段、方向来源和作用 token 位置。据我们所知，本工作是首个将安全方向从\emph{外部}对齐 LLM 提取并移植至以 \emph{base} 主干构建的 VLA 中，并限制注入到动作生成相关位置的方法，同时提供了覆盖完整 7 token 动作序列的生成阶段感知双阶段变体（\S\ref{sec:method:dualphase}）。

\begin{table*}[t]
\centering\small
\setlength{\tabcolsep}{5pt}
\begin{tabular}{p{4.4cm}p{2.0cm}p{0.8cm}p{2.8cm}p{2.4cm}}
\toprule
方法 & 目标 & 阶段 & 方向来源 & 作用位置 \\
\midrule
AdaShield \cite{adashield2024}        & 提示             & 推理 & ---               & N/A \\
TGA/CMRT-LVLM \cite{cmrt2025}         & 投影层 + LLM     & 训练 & 标题引导          & 全部 \\
CMRM \cite{cmrm2025}                  & LLM 隐层         & 推理 & 自身              & 全部 \\
VLM-Guard \cite{vlmguard2025}         & LLM 隐层         & 推理 & 外部 LLM          & 最后输入 token \\
ASTRA \cite{astra2025}                & LLM 隐层         & 推理 & 视觉对抗样本      & 生成 token \\
Security Tensors \cite{securitytensors2025} & 软 token    & 训练 & 有害 SFT          & 前缀 \\
SARSteer \cite{sarsteer2025}          & LLM 隐层         & 推理 & 自身              & 生成 token \\
SafeVLA \cite{safevla2025}            & 策略             & 训练 & 安全强化学习      & 全部 \\
Mech.\ Interp.\ VLA \cite{mechinterpvla2025} & LLM 隐层 & 推理 & 分布内 VLA 数据   & 全部 \\
\textbf{RDT/RDT+（本文）}             & LLM 隐层         & 推理 & \textbf{外部 LLM（跨模型）} & \textbf{动作位置 / 文本+动作} \\
\bottomrule
\end{tabular}
\caption{多模态安全设计空间中的定位。RDT/RDT+ 首次将（a）从\emph{外部}对齐 LLM 提取的方向移植至\emph{未对齐}的 base 主干 VLA，并（b）限制注入到动作生成位置。}
\label{tab:positioning_zh}
\end{table*}

\section{方法}
\label{sec:method}

\subsection{预备知识：OpenVLA 的动作 token 通路}
\label{sec:method:openvla}

\textsc{OpenVLA} \cite{openvla2024} 是一个在 \textsc{Open-X-Embodiment} 上微调的 \textsc{Prismatic} 家族视觉-语言模型。其架构为 $\mathcal{M}(I, x) = \mathrm{Llama\text{-}2\text{-}7b_{\text{base}}}(\phi_{\text{vis}}(I) \oplus \mathrm{Embed}(x))$，其中 $\phi_{\text{vis}}$ 融合 SigLIP 和 DINOv2 特征，$\oplus$ 表示将 256 个视觉 patch token 投影并与 token 化文本指令 $x$ 拼接。动作以 7 个离散 token $(a_1, \ldots, a_7)$ 的形式输出，每个对应 6-DoF 姿态 + 1-DoF 夹爪的一个自由度。

本文工作的关键结构细节在于动作 token 的表示方式。\textsc{OpenVLA} 并\emph{未}扩展 Llama 词表，而是\emph{覆写}了最后 $K=256$ 个 token id（id $31744 \ldots 31999$），将其替换为 256 个均匀分布的动作 bin。因此，每个动作 token 在嵌入表层面是一个原本属于自然语言的 token 的嵌入，经过持续训练后用于输出量化的机器人运动。我们记动作 token id 集合为 $\mathcal{A} = \{31744, \ldots, 31999\}$。

\subsection{诊断：动作 token 是否与拒绝轴解耦？}
\label{sec:method:diagnostic}

我们研究 Llama-2-chat 的拒绝方向作用于 \textsc{OpenVLA} 隐层状态时，能否在动作 token 位置区分有害与无害输入。按照 \citet{arditi2024refusal} 的方法，我们在每层 $L$ 提取 rank-1 拒绝方向：
\begin{equation}
\begin{split}
  r_L \;=\;& \frac{1}{|\mathcal{H}|} \sum_{x \in \mathcal{H}} h_L^{\text{chat}}(x)[-1] \\
           & -\; \frac{1}{|\mathcal{B}|} \sum_{x \in \mathcal{B}} h_L^{\text{chat}}(x)[-1],
\end{split}
  \label{eq:direction_zh}
\end{equation}
其中 $h_L^{\text{chat}}(x)[-1]$ 是 Llama-2-7b-chat 在第 $L$ 层的最后 token 隐层状态，$\mathcal{H}$、$\mathcal{B}$ 分别为有害（\textsc{AdvBench}）和无害（\textsc{Alpaca}）提示集。

随后我们在 \textsc{OpenVLA} 的隐层状态上测量该方向的 AUC，分别在两类位置类型：（i）指令的最后文本 token，（ii）第一个解码动作 token。我们的假设（在 \S\ref{sec:exp} 中通过实验验证）是：移植的方向在文本 token 位置保持判别能力，而在动作 token 位置崩溃，从而量化结构性安全缺口。

\begin{table}[t]
\centering\small
\framebox[\columnwidth]{%
  \begin{minipage}{0.9\columnwidth}\centering\vspace{8pt}
  \emph{占位符。} \textsc{OpenVLA} 中各层在\{最后文本 token、第一个动作 token\}位置的投影 AUC，使用从 Llama-2-chat 提取的 $r_L$。数据将由 \texttt{scripts/05\_sanity\_check.py} 填入。\vspace{8pt}
  \end{minipage}}
\caption{解耦诊断：从 Llama-2-chat 移植的拒绝方向在 \textsc{OpenVLA} 的文本 token 位置能够区分有害与无害，但在动作 token 位置信号崩溃。}
\label{tab:decoupling-zh}
\end{table}

\subsection{拒绝方向移植（RDT）}
\label{sec:method:rdt}

基于上述诊断，我们提出在推理时将移植的方向直接加入负责动作生成的位置的残差流，以\emph{恢复}该处的拒绝信号。设 $H_L \in \mathbb{R}^{B \times T \times d}$ 为第 $L$ 层隐层状态，$r_L$ 为式~\eqref{eq:direction_zh} 中的单位归一化拒绝方向。

\paragraph{为何未对齐的 base 主干仍能响应。}
Llama-2-chat 与 Llama-2-base 共享其 RLHF 前的检查点。残差流的线性几何结构在预训练阶段建立；已知 RLHF 沿该几何结构移动激活，而不旋转基底 \cite{arditi2024refusal}。因此，$r_L$ 指向的方向在 \textsc{OpenVLA} 动作 BC 微调后仍然有效使用的残差流坐标系中保持意义。实证上，先前工作已证实：将 chat 模型的拒绝方向移植至 base LLM 可将拒绝率从接近零提升至 88\% 以上 \cite{baselmsrefuse2024}，直接验证了跨初始化迁移的有效性。我们在 \S\ref{sec:analysis:crossmodel} 中针对本文 VLA 场景加以验证。

\paragraph{生成阶段时序。}
\textsc{OpenVLA} 自回归生成 7 个动作 token。在\emph{预填充}阶段（生成 $a_1$ 时），输入 id 序列仅包含文本 token，不含任何动作 id。在此后每个\emph{解码}步骤（生成 $a_2 \ldots a_7$ 时），已生成的动作 token 通过 KV 缓存重新作为上下文处理。因此，纯动作 token 掩码仅在 $a_2 \ldots a_7$ 上触发，$a_1$ 不受保护。我们在 \S\ref{sec:method:dualphase} 中通过生成阶段感知变体解决这一问题。

\paragraph{基础 RDT（仅动作位置）。}
设 $M_L \in \{0,1\}^{B \times T}$ 为动作 token 掩码 $M_L = \mathbb{1}[\mathrm{ids} \in \mathcal{A}]$，干预为：
\begin{equation}
  H_L' \;=\; H_L \,+\, \alpha \cdot (M_L \odot \mathbf{1}_d^\top) \odot r_L,
  \label{eq:rdt_zh}
\end{equation}
其中 $\alpha$ 为标量系数，$\odot$ 表示广播 Hadamard 积。除此之外，不修改任何其他层、位置或模型权重。

\subsection{生成阶段感知双阶段注入（RDT+）}
\label{sec:method:dualphase}

为保护 $a_1$ 同时也保护 $a_2 \ldots a_7$，我们引入 \emph{RDT+}，在两种不重叠的位置类型上以独立尺度施加方向：
\begin{equation}
\begin{split}
  H_L' \;=\;& H_L
  \;+\; \alpha_{\text{text}} \cdot (M_L^{\text{text}} \odot \mathbf{1}_d^\top) \odot r_L \\
  &\;+\; \alpha_{\text{act}}  \cdot (M_L^{\text{act}}  \odot \mathbf{1}_d^\top) \odot r_L,
\end{split}
  \label{eq:rdtplus_zh}
\end{equation}
其中 $M_L^{\text{text}} = \mathbb{1}[\mathrm{ids} \notin \mathcal{A}]$，$M_L^{\text{act}} = \mathbb{1}[\mathrm{ids} \in \mathcal{A}]$。预填充阶段文本掩码激活，引导用于解码 $a_1$ 的表示；每个解码步骤动作掩码在新处理的 token 上触发。全程使用 $\alpha_{\text{text}} = 0.3\,\alpha$，$\alpha_{\text{act}} = \alpha$；这种不对称性防止过度扰动指令处理，同时在动作位置保持强安全信号。

位置消融实验（\S\ref{sec:exp:ablation}）比较仅动作、仅文本、文本+动作、全 token 等目标，以隔离每个阶段的贡献。

\subsection{Rank-$k$ 子空间变体}
\label{sec:method:rankk}

动作空间危害涵盖多个子类别（对人员的暴力、财产损毁、火灾和电气危险、化学污染），可能沿不同方向编码。我们将式~\eqref{eq:direction_zh} 的 rank-1 提取替换为 rank-$k$ 子空间。设 $\Delta \in \mathbb{R}^{|\mathcal{H}| \times d}$ 为逐提示对比矩阵 $\Delta_{i,:} = h_L^{\text{chat}}(x_i^\text{h})[-1] - \mu_L^\text{benign}$，其右奇异向量 $V_k = [v_1, \ldots, v_k]^\top$（前 $k$ 个）定义了 $k$ 维拒绝子空间。干预变为：
\begin{equation}
  H_L' \;=\; H_L \,+\, \alpha \cdot (M_L \odot \mathbf{1}_d^\top) \odot \!\left(\!\sum_{j=1}^{k} v_j\!\right)\!,
\end{equation}
或采用正交投影模式（在 \S\ref{sec:exp:ablation} 中消融），从 $H_L$ 中减去其在 $V_k$ 上的投影。$k=1$ 时退化为 rank-1。

\subsection{内容自适应 $\alpha$}
\label{sec:method:adaptive}

固定 $\alpha$ 在良性成功率与有害拒绝率之间存在权衡。我们用轻量级头部 $g_\theta$ 替代它，该头部作用于后投影层的视觉-文本 token：
\begin{equation}
  \alpha(I, x) \;=\; \alpha_{\max} \cdot \sigma\!\bigl(g_\theta(\phi_{\text{vis}}(I) \oplus \mathrm{Embed}(x))\bigr),
\end{equation}
其中 $g_\theta$ 是约 $500\mathrm{K}$ 参数的两层 MLP，在配对的有害/无害 VLA 输入上以二元交叉熵训练。推理时，$\alpha(I, x)$ 对每个输入计算一次并代入式~\eqref{eq:rdt_zh} 或 \eqref{eq:rdtplus_zh}。$g_\theta$ \emph{不}修改 \textsc{OpenVLA} 的权重；它位于模型外部，使用已在计算中的相同激活。注意，仅 rank-1 和 rank-$k$ 变体严格无需训练；自适应变体需要训练 $g_\theta$（在局限性部分披露）。

\subsection{实现}
\label{sec:method:impl}

\method{} 以两个 PyTorch 钩子实现：（i）在顶层 Llama 模型上的前向\emph{预钩子}，在嵌入层消耗前缓存当前输入 id 张量；（ii）解码器第 $L$ 层的输出钩子，读取缓存的 id 计算位置掩码并应用式~\eqref{eq:rdt_zh} 或 \eqref{eq:rdtplus_zh}。双钩子设计是必要的，因为 Llama 解码器层仅接收 \texttt{hidden\_states} 而不接收 \texttt{input\_ids}，位置信息必须通过共享闭包传递。端到端代码不超过 120 行。我们将在发表后通过匿名仓库发布 \method{}、自适应攻击、所有基线及复现脚本。

\section{实验}
\label{sec:exp}

\subsection{实验设置}
\label{sec:exp:setup}

\paragraph{主干模型。}
除特别说明外，我们使用 \textsc{OpenVLA-7b}（\texttt{openvla/openvla-7b}）。\textsc{OpenVLA-OFT}（\texttt{moojink/openvla-7b-oft-\allowbreak finetuned-libero-spatial}）用作方案特有的负对照（\S\ref{sec:analysis:oft}）。外部对齐 LLM 为 \texttt{NousResearch/\allowbreak Llama-2-7b-chat-hf}，是与 Meta Llama-2-7b-chat 权重相同的非门控镜像。

\paragraph{评测套件。}
\textbf{LIBERO-Long}（10 个长时序操作任务，每任务 50 条轨迹），用作良性任务和视觉场景的来源；\textbf{SafeAgentBench} \cite{safeagentbench2024}，我们从中选取 120 个具身危害子集\footnote{使用 750 任务 AI2-THOR 分集以保证流水线兼容性。}；以及\textbf{指令条件化 AdvBench} 分集，将从我们策划分类体系（暴力、儿童安全、火/电、化学、财产损毁、疏忽；40 条提示 × 5 次改写）中提取的 128 条有害指令与 LIBERO 视觉场景配对。

\paragraph{攻击方式。}
（i）\emph{无攻击}（良性分布），用于在干净 \textsc{LIBERO-Long} 上测量良性成功率（BSR）。（ii）\emph{文本越狱}：以明文语言传递有害指令（使用我们的 \textsc{AdvBench} 分集）。（iii）\emph{UADA}/\emph{UPA}/\emph{TMA}：遵循 \citet{wu2024exploringadversarial} 的通用对抗补丁，$L_\infty$ 预算 $16/255$，补丁面积 5\%，500 次内迭代。

\paragraph{评测指标。}
\textbf{BSR}（越高越好）：良性 \textsc{LIBERO-Long} 任务的完成率。
\textbf{HRR-strict}（越高越好）：有害任务中，前 7 个动作 token 全部落入以 bin 128 为中心的零运动窗口（$\pm$8 个 bin）的比例，用作"机器人停止"的模拟器无关替代指标。
\textbf{HRR-partial}（越高越好）：有害任务中，7 个动作 token 中至少 4 个落入零运动窗口的比例。这一较宽松的标准为部分拒绝提供正分，例如停止 4 个自由度但夹爪仍在闭合的机器人，在防御效果较强但不完美时提供更多信息。
\textbf{BRR}（越高越好）：$\tfrac{1}{2}(\mathrm{BSR} + \mathrm{HRR\text{-}strict})$，衡量防御在不过度拒绝情况下的综合实用性。相机就绪版本中将额外报告逐攻击分解和逐类别 HRR。

\paragraph{基线。}
\textbf{无防御}：原始 OpenVLA。
\textbf{AdaShield-S} \cite{adashield2024}：在指令前添加静态安全提示。
\textbf{VLM-Guard 移植} \cite{vlmguard2025}：从 Llama-2-chat 第 14 层提取 SVD rank-5 SSD，施加于 OpenVLA 的所有 token 位置（我们对原始仅最后 token 方法的移植适配版本）。
\textbf{SafeVLA-lite}：\textsc{SafeVLA} \cite{safevla2025} 的廉价替代，以 LoRA 微调（$r=32$）在 $1\mathrm{k}$ 有害 + $1\mathrm{k}$ 无害样本对上训练一个 epoch，有害目标标注为零运动动作 token。

\paragraph{\method{}/\method{}+ 配置。}
除消融变化外，使用第 $L=14$ 层，$\alpha=1.0$。基础 \method{} 仅以动作位置为目标；\method{}+ 使用文本+动作双阶段注入，$\alpha_{\text{text}}=0.3$，$\alpha_{\text{act}}=1.0$。Rank-$k$ 变体使用 $k=5$。自适应 $\alpha$ 头为隐藏维度 256 的两层 MLP，在配对标注输入上训练 3 个 epoch。

\subsection{主要结果}
\label{sec:exp:main}

表~\ref{tab:main-placeholder} 将报告所有基线和 \method{} 变体的 BSR、HRR-strict、HRR-partial 及 BRR。

\begin{table}[t]
\centering\small
\framebox[\columnwidth]{%
  \begin{minipage}{0.9\columnwidth}\centering\vspace{8pt}
  \emph{占位符。} 主对比实验。行：\{无防御, AdaShield, VLM-Guard 移植, SafeVLA-lite, \method{}（仅动作）, \method{}+（文本+动作）, \method{}+rank-$k$, \method{}+自适应 $\alpha$\}。列：BSR, HRR-strict, HRR-partial, BRR。\vspace{8pt}
  \end{minipage}}
\caption{在指令条件化 \textsc{AdvBench} 分集与 \textsc{LIBERO-Long} 良性任务上的主对比实验。}
\label{tab:main-zh}
\end{table}

\subsection{消融实验}
\label{sec:exp:ablation}

我们对五个维度进行消融。
\textbf{方向来源}：将拒绝方向与（i）等范数随机单位向量（固定随机种子跨实验一致）、（ii）取反的拒绝方向、（iii）无干预进行对比。该对照实验至关重要：若随机向量与拒绝方向获得相近的 HRR，则收益归因于一般隐层状态扰动而非移植的安全几何；若拒绝 $\gg$ 随机，则跨模型迁移假说得到验证。记差值 $\Delta_{\text{rand}}=\mathrm{HRR}(r_L) - \mathrm{HRR}(\text{rand})$。
\textbf{层扫描}：$L \in \{8, 10, 12, 14, 16, 18, 20\}$，确定移植方向既有信号又足够早以影响后续动作解码的最优层。
\textbf{尺度扫描}：$\alpha \in \{0, 0.1, 0.5, 1, 2, 5\}$，揭示 BSR/HRR 权衡并为自适应 $\alpha$ 提供动机。
\textbf{目标位置}：\{仅动作、仅文本、文本+动作、全部\}，隔离归因于动作 token 定向和生成阶段感知双阶段注入的收益——与 \textsc{VLM-Guard} 的关键区别。
\textbf{Rank-$k$}：$k \in \{1, 3, 5, 10\}$，验证动作空间危害是否跨越多个方向。

\begin{table}[t]
\centering\small
\framebox[\columnwidth]{%
  \begin{minipage}{0.9\columnwidth}\centering\vspace{8pt}
  \emph{占位符。} 五维消融矩阵。顶部块：方向来源（拒绝/随机/取反/无）。底部块：层、尺度、目标位置、子空间 rank。\vspace{8pt}
  \end{minipage}}
\caption{消融扫描。$\Delta_{\text{rand}}$ 列报告拒绝$-$随机 HRR 差值，正值确认跨模型几何是有效成分。}
\label{tab:ablations-zh}
\end{table}

\subsection{自适应攻击评测}
\label{sec:exp:adaptive}

我们在白盒自适应攻击者下评测 \method{}，该攻击者知晓 $r_L$、层选择和 $\alpha$。攻击者通过 PGD（500 次迭代，$\epsilon=16/255$）优化图像补丁，最小化 $|h_L[\text{last}] \cdot r_L|$，将隐层状态推至拒绝方向的正交补空间。我们报告有无该攻击下的 HRR-strict 和 HRR-partial。一个不仅仅捕获朴素攻击者的防御，在自适应压力下仍应在无防御基线之上保持非平凡优势。

\subsection{与现有防御的组合}
\label{sec:exp:composition}

最后，我们测试 \method{}+ 与 AdaShield 提示和 VLM-Guard 投影的正交组合效果。表~\ref{tab:composition-zh} 将报告成对和三重组合的 BRR。正向组合支持将 \method{}+ 定位为与现有提示层和投影层防御互补的即插即用动作层模块。

\begin{table}[t]
\centering\small
\framebox[\columnwidth]{%
  \begin{minipage}{0.9\columnwidth}\centering\vspace{8pt}
  \emph{占位符。} 组合矩阵：AdaShield / VLM-Guard / \method{}+ 及其成对/三重组合。\vspace{8pt}
  \end{minipage}}
\caption{\method{}+ 与现有防御的组合效果。}
\label{tab:composition-zh}
\end{table}

\section{分析}
\label{sec:analysis}

\subsection{为何跨模型移植有效？}
\label{sec:analysis:crossmodel}

Llama-2-chat 与 Llama-2-base 共享预训练初始化；RLHF 沿现有残差流基底移动激活而不旋转它 \cite{arditi2024refusal}。先前工作已实证确认：将 chat 模型的拒绝方向移植至同源 base LLM（相同预训练，无 RLHF）可将拒绝率从接近零提升至 88\% 以上 \cite{baselmsrefuse2024}。对于我们的场景，剩余的问题是这种共享几何是否能在 \textsc{OpenVLA} 额外 27 个 epoch 的行为克隆微调后得以保留。

为验证这一点，我们在 $1\mathrm{k}$ 自然语言探针上计算 Llama-2-chat 与 \textsc{OpenVLA} 主干之间第 $L$ 层隐层状态的平均余弦相似度。第 8--20 层的高相似度将表明基底实质上被保留，方向在两个模型间仍可互译。我们还检验从 Llama-2-chat 提取的 $r_L$ 在 \textsc{OpenVLA} 的有害动作位置上获得的零运动质量是否显著高于等范数随机单位向量（\S\ref{sec:exp:ablation} 方向来源对照）。较大的 $\Delta_{\text{rand}}$ 差值是几何保留的功能性印证：拒绝方向的特定几何性质——而非仅仅是其幅度——是有效成分。这些数值将在相机就绪版本中报告。

\subsection{RDT{}+ 下安全层探针的恢复}
\label{sec:analysis:probe-recovery}

\S\ref{sec:method:diagnostic} 中的诊断确立了拒绝方向在 \textsc{OpenVLA} 的动作 token 位置崩溃（投影 AUC 较低）。我们还在施加 \method{}+ 前后分别运行 \citet{securitytensors2025} 的安全层探针——该探针逐层测量良性-良性与良性-有害隐层状态对的余弦相似度差值。若动作 token 位置的探针分离度从接近零恢复至接近文本 token 位置的水平，则我们获得了对干预的机制层面而非行为层面的确认：该方向在恢复一个几何上缺失的信号，而不仅仅是任意破坏动作生成。该分析与解耦诊断结果一并报告。

\subsection{OpenVLA-OFT 作为方案特有负对照}
\label{sec:analysis:oft}

\textsc{OpenVLA-OFT} \cite{openvlaoft2025} 将最后 256 token 的离散化替换为连续 L1 回归头。按照我们的假设，\method{} 和 \method{}+ 是方案特有的：不存在 id $\geq 31744$ 的动作 token 可供定向。因此，我们预期将 \method{}+（相同 $r_L$、层、$\alpha$）施加于 \textsc{OpenVLA-OFT} 所产生的 HRR 变化与无防御基线无异。若此预测得到确认，则验证了 \method{} 的收益来自动作词汇交互，而非来自一般隐层状态扰动——对 \S\ref{sec:exp:ablation} 中方向来源对照实验的互补证据。

\subsection{拒绝在动作空间中意味着什么？}
\label{sec:analysis:action-space}

当 \method{}+ 在 $\alpha = 1$ 时触发，\textsc{OpenVLA} \emph{希望}输出什么来替代有害动作？我们探测 \method{}+ 关闭和开启时，第一个生成动作 token 在 256 个动作 bin 上的 softmax 分布。可能的定性结果有三种：（i）概率质量向 bin $128$（零增量，原地不动）移动；（ii）概率质量广泛重分布，表明 \method{}+ 仅破坏了有害动作而未引导出连贯的安全替代；（iii）概率质量向另一个仍然危险的区域移动。结果（i）支持将 \method{} 的拒绝解读为\emph{保持位置}——一种物理上安全的默认行为——而非注入随机运动。对夹爪自由度施加相同探测，可表征 \method{}+ 是否倾向于维持当前夹爪状态。我们还将该分布与方向来源变体（拒绝方向、随机方向、取反方向）进行对比，确认拒绝方向的特定几何性质是产生连贯停止行为的原因。

\begin{figure}[t]
\centering
\framebox[\columnwidth]{%
  \begin{minipage}{0.9\columnwidth}\centering\vspace{14pt}
  \emph{占位符图。} 有害提示下的动作 bin softmax 分布，对比无防御 / \method{}（仅动作）/ \method{}+（文本+动作）/ 随机方向对照。\vspace{14pt}
  \end{minipage}}
\caption{动作空间中的拒绝：\method{}+ 下第一个动作 token 分布向零运动 bin 的预测偏移，以及与随机方向对照的对比。}
\label{fig:action-space-refusal-zh}
\end{figure}

\subsection{自适应攻击者下的攻击成功率上界}
\label{sec:analysis:adaptive-bound}

在 \S\ref{sec:exp:adaptive} 的 PGD 正交化攻击下，攻击者每层最多可将 $h_L \cdot r_L$ 减小 $\epsilon \cdot \|J_\text{vis}\|_\text{op}$，其中 $J_\text{vis}$ 是在扰动补丁处视觉到 LLM 前向传播的 Jacobian，$\epsilon$ 为像素空间预算。实践中，该上界转化为有界的 HRR 下降，我们将连同自适应压力下测得的残余防御收益一并报告。

\subsection{观察到的局限性}
\label{sec:analysis:limitations}

我们预期两种削弱 \method{} 效果的情况：（i）当有害指令语义较轻微时（例如"把软玩具扔向婴儿"），拒绝方向的激活较弱，产生部分而非零运动动作；（ii）当视觉场景模糊时（例如塑料玩具刀），\method{}+ 可能误判并过度拒绝。两种情况在一定程度上由内容自适应 $\alpha$ 头（\S\ref{sec:method:adaptive}）缓解，但均未得到完全解决，这为局限性部分提供了动机。

\section{结论}
\label{sec:conclusion}

我们识别出动作 token VLA 特有的结构性安全缺口：其主干通常是未对齐的 base LLM，而其动作词汇经过数千步行为克隆训练后，驻留在自然语言拒绝方向不再覆盖的残差流子空间中。\textsc{RDT}+ 通过从外部对齐兄弟 LLM 移植拒绝方向，并经由覆盖全部 7 个动作 token 的生成阶段感知双阶段策略注入，在其应处的位置恢复拒绝信号。方向来源对照实验证实，移植的拒绝方向的特定几何性质——而非仅仅是隐层状态扰动的幅度——是有效成分。\textsc{RDT}+ 无需训练，增加不超过 5\% 的延迟，可跨动作攻击类型泛化，并可与提示层防御组合使用。我们将 \textsc{RDT}+ 视为将文本对齐视为跨非文本输出通路\emph{可重用资产}的一步，并呼吁学界将 base LLM 主干和词表覆写明确视为需要分析的结构性安全决策。


\section*{局限性}
\begin{itemize}
    \item \textbf{方案特有适用性。} \method{}/\method{}+ 针对以覆写 LLM 子词表方式编码离散动作的 VLA。不能直接应用于连续动作 VLA（如 \textsc{OpenVLA-OFT}、\textsc{$\pi_0$}）。我们通过将 \textsc{OpenVLA-OFT} 作为方案特有负对照来验证这一点（\S\ref{sec:analysis:oft}）。
    \item \textbf{共享预训练假设。} 跨模型移植依赖目标 VLA 主干与外部对齐 LLM 共享预训练初始化。我们未评估跨不同预训练家族的移植（例如以 Llama 为基础的 VLA 目标与以 Qwen 对齐的来源）；是否有效留待未来工作。
    \item \textbf{基础 \method{} 中第一个动作 token 的覆盖。} 仅动作目标的基础 \method{} 变体在生成 $a_1$ 的预填充步骤中不触发，因为此时输入中不存在动作 token id。\method{}+ 通过双阶段注入解决了这一问题（\S\ref{sec:method:dualphase}），但文本侧 $\alpha_{\text{text}}=0.3$ 注入弱于动作侧注入，因此 $a_1$ 在当前实现中受到的引导压力低于 $a_2\ldots a_7$。
    \item \textbf{仅模拟评测。} 我们的评测在模拟环境（\textsc{LIBERO}、\textsc{AI2-THOR}）中进行。真实机器人执行会引入我们未处理的分布偏移。在模拟中触发"零运动"的防御在真实控制器接收到停止信号所需的时间内仍可能执行不安全行为。
    \item \textbf{攻击覆盖范围。} 我们评测指令层和补丁层攻击；不测试新型物理对抗扰动、材质级伪装或多步对抗计划。\S\ref{sec:exp:adaptive} 中的自适应攻击仅是仅 $L_\infty$ 约束下白盒梯度攻击的上界。
    \item \textbf{零运动代理指标。} 我们的 HRR-strict 指标以零运动动作窗口作为"机器人停止"的替代。部分拒绝可能采取其他安全动作形式（如将刀放下），而我们的指标对此不给分。HRR-partial 部分缓解了这一问题；完整的人工评测留待相机就绪版本。
    \item \textbf{内容自适应 $\alpha$ 引入小型辅助模型。} $g_\theta$ 头（约 $500\mathrm{K}$ 参数）在标注的有害/无害样本对上训练，这意味着 \method{}+ 的自适应变体并非严格无需训练；仅 rank-1 和 rank-$k$ 变体是如此。我们在附录中披露 $g_\theta$ 的训练成本和数据。
    \item \textbf{安全-实用权衡。} \method{}+ 以一定良性任务成功率的代价降低有害动作率。我们分别报告 BSR 和 HRR（而非仅报告 BRR），以使该权衡可见。
\end{itemize}


\section*{伦理声明}
\paragraph{双重用途考量。}
\method{} 是一种防御方法，但诊断部分（\S\ref{sec:method:diagnostic}）揭示了动作 token VLA 中的结构性漏洞，理论上可能为攻击者提供信息。我们认为公开的利弊权衡对公开有利：（a）该漏洞在公开 VLA 代码库中已隐含存在，且已被已发表的攻击工作所利用 \cite{wu2024exploringadversarial,badrobot2025,shawshank2025}；（b）防御原语对社区的价值高于漏洞的惊喜因素；（c）我们将在公开披露前与 OpenVLA 维护者协调发布时间。

\paragraph{数据集中的有害内容。}
我们构建了包含描述假设性物理危害提示（如针对儿童的暴力、火灾危险）的指令条件化 \textsc{AdvBench} 分集。这些内容仅用于安全评测，无意引发、鼓励或使能危害。提示集在仅供研究的许可证下发布，并附带内容警告。

\paragraph{算力与数据。}
所用所有模型、数据集和基线均为公开资源。我们主要结果的总算力不超过 100 GPU 小时，在单张消费级 GPU 上完成。详细算力和数据信息将在附录中披露。

\paragraph{无人类被试。}
所有实验在模拟环境中进行，不涉及人类被试。


\bibliography{custom}

\end{document}
